\chapter*{About the Author}
\addcontentsline{toc}{chapter}{About the Author}
\thispagestyle{plain}

\textbf{Abhijeet Kumar} is a mechanical engineer and computational researcher with a Ph.D.\ from IIT Bombay. His work sits at the intersection of computational mechanics, finite element analysis, structural engineering, scientific machine learning and engineering automation.

His experience covers finite element analysis, structural dynamics, vibration and shock analysis, nonlinear numerical modeling, reduced-order and surrogate modeling, structural health monitoring, signal processing and engineering data analytics. He is particularly interested in combining physics-based simulation with machine learning to build faster, more reliable engineering tools.

His doctoral research at IIT Bombay produced machine-learning frameworks for assessing debonding in stiffened structures. They combine finite element simulations, experimental vibration data, signal processing, feature extraction, deep learning, autoencoders and meta-learning, with a focus on closing the gap between simulated and experimental data and on staying robust when the data conditions change.

His current interests include physics--data hybrid modeling, reduced-order and surrogate modeling, differentiable simulation, physics-aware machine learning, system identification, uncertainty quantification, sensitivity analysis and engineering optimization. He develops Python frameworks that integrate finite element simulation, PyTorch and GPU computation, model reduction, calibration and uncertainty-aware analysis. This book grew out of that work: it is the written account of two packages, \fea{} and \rom{}, built from scratch so that every step from mesh to reduced model can be read, run and checked.

\medskip
\noindent\textbf{Profile:} \url{https://www.linkedin.com/in/abhijeet-kumar-phd-a689b429}
